В статье рассматривается задача автоматизированной оцифровки химической литературы с использованием методов машинного зрения и анализа изображений. Актуальность работы обусловлена тем, что значительная часть накопленных химических знаний хранится в сканированных книгах, статьях и архивных документах, где наряду с текстом присутствуют структурные формулы, схемы реакций, таблицы и иные графические элементы; их автоматическая оцифровка требует специализированных методов обработки, поскольку стандартное распознавание текста не позволяет извлечь смысл графических химических объектов.
Для решения этой проблемы предложен конвейер обработки химических документов, включающий адаптивную предобработку изображений, сегментацию страниц, выделение текстовых и графических областей, распознавание текста и химических структур, а также постобработку и преобразование результатов в машиночитаемый формат SMILES.
Экспериментальная апробация проведена на изображениях страниц химических книг. Для распознавания текста лучшее качество (среднее значение метрики $1-\mathtt{WER}$ равное 0,7693) показала комбинированная схема, сочетающая адаптивную предобработку изображения с последующей языковой нормализацией текста, что в среднем повышает качество работы OCR примерно на 35\% по сравнению с базовой конфигурацией. Для распознавания химических структур сравнивались модели MolScribe, DECIMER и Qwen3VL: адаптивная предобработка изображения не дала универсального прироста качества --- лучший результат без предобработки показала MolScribe, тогда как после предобработки её качество снизилось, результат DECIMER не изменился, а Qwen3VL продемонстрировала лишь незначительное улучшение.
Показано, что эффективность предобработки изображений зависит от конкретной модели и типа химической нотации, поэтому её применение требует индивидуальной настройки под задачу, а не может использоваться как универсальный этап обработки.
